\chapter{Catálogo de indicadores de compromiso (IoCs)}

En este apéndice se compendian los artefactos maliciosos, rutas sospechosas y atributos de persistencia identificados durante la investigación pericial en el servidor \texttt{helios}.

\section{Rutas y binarios clave identificados}

\begin{table}[h!]
\centering
\footnotesize
\begin{tabularx}{\textwidth}{@{} l l X @{}}
\toprule
\textbf{Ruta del artefacto} & \textbf{Propietario / permisos} & \textbf{Descripción / función detectada} \\
\midrule
\texttt{/dev/ptyp} & \texttt{root:root (drwxr-xr-x)} & Directorio oculto de depósito y scripts \\
\texttt{/dev/ptyp/energyb} & \texttt{root:root (rwxr-xr-x)} & Bot IRC con soporte DCC y retransmisión \\
\texttt{/dev/ptyp/psybnc} & \texttt{root:root (rwxr-xr-x)} & Bouncer IRC para encubrimiento de conexiones \\
\texttt{/usr/lib/popclean} & \texttt{root:root (-u--ia-----)} & Binario modificado con atributo inmutable \\
\texttt{/usr/local/lib/dsniff.services} & \texttt{root:root (-u--ia-----)} & Capturador y analizador de credenciales de red \\
\texttt{/tmp/...} & \texttt{duende:duende (drwxrwxrwx)} & Directorio temporal de desempaque de malware \\
\bottomrule
\end{tabularx}
\caption{Artefactos y ejecutables sospechosos identificados en el sistema comprometido.}
\end{table}

\section{Atributos especiales de archivos del sistema}

Varios de los ejecutables troyanizados fueron fijados con atributos especiales de Ext2 para evitar su eliminación o modificación incluso por el usuario \texttt{root}:
\begin{itemize}
    \item \textbf{Atributo inmutable (\texttt{+i}):} Impide que el archivo sea modificado, renombrado, eliminado o sobreescrito. Solo puede removerse restableciendo explícitamente el atributo con el comando \texttt{chattr -i}.
    \item \textbf{Atributo de adición únicamente (\texttt{+a}):} Restringe la apertura del archivo únicamente para agregar información al final, bloqueando el truncamiento o borrado.
\end{itemize}
